\subsection{Related work}
\paragraph{Adaptation, retention and parameterization.}
Within the LoRA family, several methods use singular directions of pretrained weights or gradients to initialize adaptation \citep{meng2024pissa,wang2024loraga,zhang2025loraone,ponkshe2024lorasb}. 
Optimization and scaling studies explain how parameterization affects adaptation \citep{hayou2024loraplus,kalajdzievski2023rslora,wang2024loraga}; \citet{lee2026learningrate} link preferred learning rates to loss curvature. \citet{biderman2024loraforgetting} compare LoRA's adaptation and forgetting with full fine-tuning. PEFT-Arena compares PEFT methods using a shared SFT learning rate and one run per configuration, leaving method-specific tuning and seed uncertainty unresolved \citep{huang2026peftarena}. ETHER questions the practical necessity of hyperspherical-energy preservation \citep{bini2024ether}. We examine preservation alongside tuned retention comparisons and spectral interventions.

\paragraph{Spectral structure and functional interventions.}
\citet{shuttleworth2024illusion} identify intruder singular directions after LoRA adaptation. \citet{he2025posttrainingstructure} study spectrum and vector replacement after full fine-tuning. \citet{jin2026rlheals} connect forgetting and recovery during full-parameter SFT/RL to singular-vector changes. CASA studies spectral preservation in diffusion-LoRA transfer \citep{wang2026casa}, while SVFT learns combinations of pretrained singular-vector pairs \citep{lingam2024svft}. Our contribution is to connect tuning-controlled comparisons with complementary keep/reset interventions, distinguishing recurring retention responses from setting-dependent contributions to task performance.
